\subsection{A renormalisation-group recursion for the cube constant}\label{rR6:sec}

This subsection gives an exact form of the upward step of
\cite[Sec.~X\,K]{BBLR26}. There \cite[Theorem~100]{BBLR26} passes
$\mathrm{BP}$ of Definition~\ref{def:avgcube} from the law of the block
variables to the fine measure under the split inequality $\mathrm S$, a
sufficient condition. Here the fine constant is characterised exactly, for
every probability measure and every sub-$\sigma$-algebra of block variables
with a positive fibre constant, by a one-parameter family of effective
quadratic forms on the block variables, the Feshbach--Schur complements of
the fine Dirichlet form, as in the spectral renormalisation group
of~\cite{BFS98}. Comparing these forms with the coarse Dirichlet form gives
a two-sided recursion for the best constants along a tower. The comparison
involves three quantities of one block-spin step: the fibre constant, the
renormalisation ratio and the fluctuation response. For
translation-invariant measures on scalar fields and block sums the
recursion is an identity on the constant linear function, with the defect
$\delta_j=1-\Var(X\mid\text{fine exterior})/\Var(X\mid\text{coarse
exterior})$ of the mass $X$ of one cube. For the free massive scalar field
this defect is at most $\Gamma_d(s)/n_{j-1}$, so the defects are summable
over the scales of a tower of fixed physical cube and their product lies
between $(1+\Gamma_d(s)/n_0)^{-1}$ and $1$; in this factorisation the
factor $s^2$ of Theorem~\ref{ffs:thm:exact} enters through the coarsest
scale and the susceptibility $1/\mu^2$. For $\SU(2)$ the subsection says
which of Ba{\l}aban's estimates bear on each quantity and states, as the
open lemma (L) of Definition~\ref{rR6:lemmaL}, a condition across scales
that, with (H1), would give $\mathrm{BP}$ along the tuning. That condition
is not proved, and neither (H1) nor a mass gap is proved here.

\emph{Notation from~\cite{BBLR26}.} A block-spin map of range $r$ with
blocks of side $b$ is a measurable $c:G^{E'}\to G^E$ from the links $E'$ of
$\Lambda_{bL}$ to the links $E$ of $\Lambda_L$ such that $c_{(y,\nu)}$
depends only on the links based in the blocks $B_w=bw+[0,b)^4$ with
$\lvert w-y\rvert_\infty\le r$, and every gauge transformation $g$ of
$\Lambda_{bL}$ satisfies $c\circ g=\hat g\circ c$ for a gauge transformation
$\hat g$ of $\Lambda_L$, every one of which arises; the straight-path map
is an example. For a gauge-invariant probability measure $\mu'$ on
$G^{E'}$, $\mu^c=c_*\mu'$ is gauge invariant; $\mathcal A'$, $\mathcal A^c$
are the gauge-invariant functions in $L^2(\mu')$, $L^2(\mu^c)$, and
$\mathcal E'_{n'}(f)=n'^{-4}\sum_{x'}\mu'((f-P'_{x'}f)^2)$, $\mathcal
E^c_n(F)=n^{-4}\sum_y\mu^c((F-P^c_yF)^2)$ are the Dirichlet forms of the
cube chains of Definition~\ref{def:avgcube}. For real $f\in\mathcal A'$
there is a real $F_f\in\mathcal A^c$ with $\E[f\mid c]=F_f\circ c$; then
$\Var_{\mu'}(f)=\mu'(\Var(f\mid c))+\Var_{\mu^c}(F_f)$ and
$\mu^c((F_f-P^c_yF_f)^2)=\mu'(\Var(\E[f\mid c]\mid\mathcal G_y))$, where
$\mathcal G_y\subset\sigma(c)$ is generated by the $c_{(z,\nu)}$ with
$z\notin y+[0,n)^4$ \cite[Lemma~98]{BBLR26}; and $F\circ c\in\mathcal A'$
for real $F\in\mathcal A^c$. The \emph{split inequality}
$\mathrm S(n',n;C_A,C_B)$ for $(\mu',c)$ is
\[
\frac{\mu'(\Var(f\mid c))}{C_A}+\frac{\mathcal E^c_n(F_f)}{C_B}\le
\mathcal E'_{n'}(f)\quad\text{for every real }f\in\mathcal A',
\]
and the \emph{separate inequalities} are \textup{(A)}
$\mu'(\Var(f\mid c))\le C_A\mathcal E'_{n'}(f)$ and \textup{(B)} $\mathcal
E^c_n(F_f)\le C_B\mathcal E'_{n'}(f)$ for every real $f\in\mathcal A'$
\cite[Definition~99]{BBLR26}. By \cite[Theorem~100]{BBLR26}, $\mathrm S$
and $\mathrm{BP}_{\mu^c}(n,\varpi_c)$ give $\mathrm{BP}_{\mu'}(n',\min(1/
C_A,\varpi_c/C_B))$, and (A), (B) give $\mathrm{BP}_{\mu'}(n',(C_A+C_B/
\varpi_c)^{-1})$. A tower is a sequence of block-spin maps $c_j$ from
$\Lambda_{L_j}$ to $\Lambda_{L_{j-1}}$, $L_j=b^jL_0$, $1\le j\le K$, with
measures $\mu_{j-1}=(c_j)_*\mu_j$ and cube sides $n_j$, $L_j\ge2n_j+4$;
\cite[Theorem~101]{BBLR26} iterates $\mathrm S$ along it. Along the tuning
\cite[Corollary~102]{BBLR26} takes $a_k=2^{-k}a_0$, $n_k=2^kn_0$,
$\ell=n_ka_k$, tori $\Lambda_L$ with $L=2^kL_0$, $L_0$ even and $L_0\ge
2n_0+4$, in an infinite set $\mathcal L_k$, the Wilson measure $\mu_{k,L}$
at $\beta_k$, maps $c_1,\dots,c_k$ with blocks of side two, and the tower
$\mu^{(j)}_{k,L}$ with $\mu^{(k)}_{k,L}=\mu_{k,L}$; its hypotheses are
\textup{(H1)} $\mathrm{BP}(n_0,\varpi_0)$ for $\mu^{(0)}_{k,L}$, and
\textup{(H2)} $\mathrm S(n_j,n_{j-1};C_{A,j},C_{B,j})$ for
$(\mu^{(j)}_{k,L},c_j)$ with $C_{A,j}\le C_A$ and $\prod_{i=j}^kC_{B,i}\le
\Gamma$, uniformly in $k$ and $L\in\mathcal L_k$, and its conclusion is
$\mathrm{BP}(n_k,\min(\varpi_0,1/C_A)/\Gamma)$ for $\mu_{k,L}$. Finally
\cite[Theorem~94(i)]{BBLR26}: if $n\ge2r+2$, $L\ge2n+4$ is even, $2\le n'\le
b(n-2r-1)+1$ and $bL\ge2n'+4$, then $\mathcal E'_{n'}(F\circ c)\le(bn/n')^4
\mathcal E^c_n(F)$ for every real $F\in\mathcal A^c$.

\subsubsection{The exact step}

\emph{Setting.} Let $\mathcal H$ be a real Hilbert space, $A$ a bounded
self-adjoint operator on $\mathcal H$ with $A\ge0$, and $Q$ an orthogonal
projection on $\mathcal H$. Put $\mathcal H_c=Q\mathcal H$, $\mathcal H_f=
(I-Q)\mathcal H$, and write $A_{cc}=QAQ$, $A_{fc}=(I-Q)AQ$, $A_{cf}=
A_{fc}^{\mathsf T}$ and $A_{ff}=(I-Q)A(I-Q)$, regarded as operators between
$\mathcal H_c$ and $\mathcal H_f$. Let $A^c\ge0$ be a bounded self-adjoint
operator on $\mathcal H_c$. Put
\begin{gather*}
\varpi'=\inf_{f\ne0}\frac{\ip f{Af}}{\lVert f\rVert^2},\qquad
\varpi_c=\inf_{0\ne g\in\mathcal H_c}\frac{\ip g{A^cg}}{\lVert g\rVert^2},
\\
\varpi_{\rm fib}=\inf_{0\ne k\in\mathcal H_f}\frac{\ip k{Ak}}{\lVert k
\rVert^2},
\end{gather*}
with $\varpi_{\rm fib}=+\infty$ if $\mathcal H_f=\{0\}$. Restricting the
infimum defining $\varpi'$ to $\mathcal H_f$ gives $\varpi'\le\varpi_{\rm
fib}$, so $\varpi'=0$ if $\varpi_{\rm fib}=0$; assume $\varpi_{\rm fib}>0$.
For $0\le\lambda<\varpi_{\rm fib}$ and $g\in\mathcal H_c$ let
\[
\tilde{\mathcal E}_\lambda(g)=\inf_{k\in\mathcal H_f}\bigl(\ip{g+k}{A(g+k)}
-\lambda\lVert k\rVert^2\bigr),
\]
the effective form at the spectral parameter $\lambda$. Let $k_0(g)$ be the
minimiser at $\lambda=0$ (Theorem~\ref{rR6:feshbach}(i)), the fluctuation
response of $g$.

\emph{The measure-theoretic case.} Let $(\Omega,\mathfrak F,\mu)$ be a
probability space, $(\mathfrak F_x)_{x\in X'}$ a finite family of
sub-$\sigma$-algebras with conditional expectations $\Pi_x$, $w'>0$, and
$\mathfrak B\subset\mathfrak F$ a sub-$\sigma$-algebra (the block variables)
with conditional expectation $Q$. Let $(\mathfrak G_y)_{y\in X}$ be a finite
family of sub-$\sigma$-algebras of $\mathfrak B$ with conditional
expectations $\Pi_{\mathfrak G_y}$, and $w>0$. Let $\mathcal H$ be a closed
subspace of the real functions in $L^2(\mu)$ with mean zero that is mapped
into itself by every $\Pi_x$ and by $Q$, and suppose that every
$\Pi_{\mathfrak G_y}$ maps $\mathcal H_c=Q\mathcal H$ into itself. Take
\[
A=w'\sum_{x\in X'}(I-\Pi_x),\qquad A^c=w\sum_{y\in X}(Q-\Pi_{\mathfrak G_y})
\]
on $\mathcal H$ and on $\mathcal H_c$. Since $\mathfrak G_y\subset\mathfrak
B$, $\Pi_{\mathfrak G_y}Q=Q\Pi_{\mathfrak G_y}=\Pi_{\mathfrak G_y}$, so each
$Q-\Pi_{\mathfrak G_y}$ is an orthogonal projection, $A^c$ is $w$ times a
sum of orthogonal projections, and $\ip g{A^cg}=w\sum_y\lVert g-
\Pi_{\mathfrak G_y}g\rVert^2$ for $g\in\mathcal H_c$. This contains the
setting of \cite[Sec.~X\,K]{BBLR26}: take $\mu=\mu'$, $\mathfrak F_{x'}$
the links outside the fine cube $V'_{x'}$, $w'=n'^{-4}$, $\mathfrak B=
\sigma(c)$, $\mathfrak G_y$ as in the notation above, $w=n^{-4}$ and
$\mathcal H$ the real functions in $\mathcal A'$ of mean zero. Then $\ip f
{Af}=\mathcal E'_{n'}(f)$; by \cite[Lemma~98]{BBLR26}, its proof and gauge covariance,
$\mathcal H_c$ is the set of functions $F\circ c$ with real $F\in\mathcal
A^c$ of mean zero, $\Pi_{\mathfrak G_y}(F\circ c)=(P^c_yF)\circ c$, $\ip
{F\circ c}{A^c\,F\circ c}=\mathcal E^c_n(F)$ and $\lVert k\rVert^2=\mu'(\Var
(f\mid c))$ for $f=F\circ c+k$, $k\in\mathcal H_f$; so $\varpi'$ and
$\varpi_c$ are the best constants of $\mathrm{BP}_{\mu'}(n')$ and
$\mathrm{BP}_{\mu^c}(n)$. Product spaces with a block-averaging map are the
case $\Omega=S^{\Lambda'}$, $\mathfrak B=\sigma(c)$ for a measurable $c:
S^{\Lambda'}\to T^{\Lambda}$ whose coordinate $c_y$ depends on the spins of
the block $B_y$, with $\mathfrak F_{x'}$ and $\mathfrak G_y$ the spins
outside a fine cube and the block variables outside a coarse cube.
\begin{theorem}[the exact step]\label{rR6:feshbach}
In the setting above:
\begin{enumerate}[label=\textup{(\roman*)},leftmargin=*]
\item For $0\le\lambda<\varpi_{\rm fib}$ and $g\in\mathcal H_c$ the infimum
defining $\tilde{\mathcal E}_\lambda(g)$ is attained exactly at
$k_\lambda(g)=-(A_{ff}-\lambda)^{-1}A_{fc}g$, and
\[
\tilde{\mathcal E}_\lambda(g)=\ip g{A_{cc}g}-\ip{A_{fc}g}{(A_{ff}-
\lambda)^{-1}A_{fc}g}.
\]
The map $\lambda\mapsto\tilde{\mathcal E}_\lambda(g)$ is nonincreasing and
concave, and $\tilde{\mathcal E}_0(g)=\min\{\ip fAf:Qf=g\}$.
\item For $0\le\lambda<\varpi_{\rm fib}$, $\varpi'\ge\lambda$ holds if and
only if $\tilde{\mathcal E}_\lambda(g)\ge\lambda\lVert g\rVert^2$ for every
$g\in\mathcal H_c$. With $s=\sup\{\lambda\in[0,\varpi_{\rm fib}):
\tilde{\mathcal E}_\lambda\ge\lambda\text{ on }\mathcal H_c\}$,
\[
\varpi'=\min(s,\varpi_{\rm fib}).
\]
\item For $0\le\lambda<\varpi_{\rm fib}$ and $g\in\mathcal H_c$,
\[
\tilde{\mathcal E}_0(g)-\frac{\lambda\,\varpi_{\rm fib}}{\varpi_{\rm fib}-
\lambda}\lVert k_0(g)\rVert^2\le\tilde{\mathcal E}_\lambda(g)\le\tilde
{\mathcal E}_0(g)-\lambda\lVert k_0(g)\rVert^2 .
\]
\end{enumerate}
\end{theorem}

\begin{proof}
(i) For $k\in\mathcal H_f$, $\ip{g+k}{A(g+k)}-\lambda\lVert k\rVert^2=\ip g
{A_{cc}g}+2\ip k{A_{fc}g}+\ip k{(A_{ff}-\lambda)k}$. Since $A_{ff}\ge
\varpi_{\rm fib}$ on $\mathcal H_f$, $R=A_{ff}-\lambda\ge\varpi_{\rm fib}-
\lambda>0$ is invertible on $\mathcal H_f$ with a bounded positive inverse,
and
\begin{multline*}
2\ip k{A_{fc}g}+\ip k{Rk}\\=\lVert R^{1/2}k+R^{-1/2}A_{fc}g\rVert^2-\ip
{A_{fc}g}{R^{-1}A_{fc}g},
\end{multline*}
which is minimal exactly when $R^{1/2}k=-R^{-1/2}A_{fc}g$, that is $k=
k_\lambda(g)$. For fixed $k$ the bracket is an affine nonincreasing
function of $\lambda$; an infimum of such functions is nonincreasing and
concave. At $\lambda=0$ the functions $f$ with $Qf=g$ are the $g+k$,
$k\in\mathcal H_f$.

(ii) If $\varpi'\ge\lambda$, then for $k\in\mathcal H_f$, since $g\perp k$,
$\ip{g+k}{A(g+k)}-\lambda\lVert k\rVert^2\ge\lambda\lVert g+k\rVert^2-
\lambda\lVert k\rVert^2=\lambda\lVert g\rVert^2$, and the infimum over $k$
gives $\tilde{\mathcal E}_\lambda(g)\ge\lambda\lVert g\rVert^2$. Conversely,
for $f=g+k$ with $g=Qf$,
\begin{align*}
\ip fAf-\lambda\lVert f\rVert^2&=\bigl(\ip fAf-\lambda\lVert k\rVert^2
\bigr)-\lambda\lVert g\rVert^2\\
&\ge\tilde{\mathcal E}_\lambda(g)-\lambda\lVert g\rVert^2\ge0 .
\end{align*}
The set of admissible $\lambda$ contains $0$, since $\tilde{\mathcal E}_0
\ge0$, and is an interval, since $\tilde{\mathcal E}_\lambda(g)-\lambda
\lVert g\rVert^2$ is nonincreasing in $\lambda$. Restricting the infimum
defining $\varpi'$ to $\mathcal H_f$ gives $\varpi'\le\varpi_{\rm fib}$. For
every admissible $\lambda$, $\varpi'\ge\lambda$, so $\varpi'\ge s$. If $s<
\varpi_{\rm fib}$, every $\lambda\in(s,\varpi_{\rm fib})$ is not admissible,
so $\varpi'<\lambda$, and $\varpi'\le s$. If $s=\varpi_{\rm fib}$, then
$\varpi'\ge s=\varpi_{\rm fib}\ge\varpi'$.

(iii) Put $a=A_{fc}g$, so that $k_0(g)=-A_{ff}^{-1}a$. By (i) and the
resolvent identity $R^{-1}-A_{ff}^{-1}=\lambda A_{ff}^{-1}R^{-1}$, in which
all operators are functions of $A_{ff}$ and commute,
\[
\tilde{\mathcal E}_0(g)-\tilde{\mathcal E}_\lambda(g)=\lambda\ip a{A_{ff}
^{-1}R^{-1}a}=\lambda\ip{k_0(g)}{\phi(A_{ff})k_0(g)},
\]
with $\phi(t)=t/(t-\lambda)$. On the spectrum of $A_{ff}$, contained in
$[\varpi_{\rm fib},\infty)$, $1\le\phi(t)\le\varpi_{\rm fib}/(\varpi_{\rm
fib}-\lambda)$, so $\lVert k_0\rVert^2\le\ip{k_0}{\phi(A_{ff})k_0}\le
\varpi_{\rm fib}\lVert k_0\rVert^2/(\varpi_{\rm fib}-\lambda)$.
\end{proof}

In the measure-theoretic case $\tilde{\mathcal E}_0(g)$ is the least fine
Dirichlet form of a function whose conditional mean given the block
variables is $g$, and $k_0(g)$ is the part of that optimal function which
has conditional mean zero. Part (ii) is an equivalence at each $\lambda$, so
every sufficient condition for the upward step that gives $\varpi'\ge
\lambda$ with $\lambda<\varpi_{\rm fib}$ amounts to the lower bound
$\tilde{\mathcal E}_\lambda\ge\lambda$ on $\mathcal H_c$. The split
inequality is one such condition.

\begin{proposition}[the split inequality as a special case]\label{rR6:split}
If $\mathrm S(n',n;C_A,C_B)$ holds for $(\mu',c)$, then $\varpi_{\rm fib}
\ge1/C_A$ and $\tilde{\mathcal E}_\lambda(F\circ c)\ge\mathcal E^c_n(F)/
C_B$ for $0\le\lambda<1/C_A$ and every real $F\in\mathcal A^c$ of mean zero;
so Theorem~\textup{\ref{rR6:feshbach}(ii)} gives $\varpi'\ge\min(1/C_A,
\varpi_c/C_B)$, which is \cite[Theorem~100(i)]{BBLR26}.
\end{proposition}

\begin{proof}
For $k\in\mathcal H_f$, $\E[k\mid c]=0$, so $F_k$ is constant and
$\mathrm S$ with $f=k$ gives $\ip kAk\ge\lVert k\rVert^2/C_A$. For $f=
F\circ c+k$ one may take $F_f=F$, and $\mathrm S$ gives
$\ip fAf-\lambda\lVert k\rVert^2\ge(1/C_A-\lambda)\lVert k\rVert^2+
\mathcal E^c_n(F)/C_B\ge\mathcal E^c_n(F)/C_B$. If $\lambda<\min(1/C_A,
\varpi_c/C_B)$, then $\tilde{\mathcal E}_\lambda(F\circ c)\ge\varpi_c
\lVert F\rVert^2/C_B\ge\lambda\lVert F\circ c\rVert^2$, so $\varpi'\ge
\lambda$; let $\lambda$ increase to the minimum.
\end{proof}
\subsubsection{The recursion}

Three quantities of one step compare the effective forms with the coarse
form. With the conventions $0/0=0$ and that the infimum and supremum run
over $g\in\mathcal H_c$ with $\ip g{A^cg}>0$, put
\begin{gather*}
\kappa=\inf_g\frac{\tilde{\mathcal E}_0(g)}{\ip g{A^cg}},\qquad
\mathrm K=\sup_g\frac{\tilde{\mathcal E}_0(g)}{\ip g{A^cg}},\\
\zeta=\sup_{0\ne g\in\mathcal H_c}\frac{\lVert k_0(g)\rVert^2}{\lVert g
\rVert^2},
\end{gather*}
the lower and upper renormalisation ratios and the fluctuation response of
the step, and assume that $\tilde{\mathcal E}_0(g)=0$ whenever $\ip g{A^cg}
=0$ when $\mathrm K$ is used. In the measure-theoretic case of
\cite[Sec.~X\,K]{BBLR26}, $\kappa\ge1/C_B$ whenever the separate
inequality (B) holds, and the best constant in (B) is $1/\kappa$, since (B)
says $\mathcal E^c_n(F_f)\le C_B\mathcal E'_{n'}(f)$ for every $f$ with
$\E[f\mid c]=F_f\circ c$, that is $\mathcal E^c_n(F)\le C_B\tilde
{\mathcal E}_0(F\circ c)$ for every real $F\in\mathcal A^c$ of mean
zero.

\begin{corollary}[the two-sided recursion]\label{rR6:recursion}
In the setting of Theorem~\textup{\ref{rR6:feshbach}},
\begin{align*}
\min\Bigl(\frac{\varpi_{\rm fib}}2,\ \frac{\kappa\,\varpi_c}{1+2\zeta}
\Bigr)&\le\varpi'\\
&\le\min\Bigl(\varpi_{\rm fib},\ \inf_{0\ne g\in\mathcal H_c}
\frac{\tilde{\mathcal E}_0(g)}{\lVert g\rVert^2+\lVert k_0(g)\rVert^2}
\Bigr)\\
&\le\min(\varpi_{\rm fib},\mathrm K\varpi_c).
\end{align*}
More precisely, $\varpi'\ge\lambda$ for every $\lambda\in[0,\varpi_{\rm
fib})$ with $\kappa\varpi_c-\lambda\zeta\varpi_{\rm fib}/(\varpi_{\rm fib}-
\lambda)\ge\lambda$. If $\zeta=0$ and $\kappa=\mathrm K$, then $\varpi'=
\min(\varpi_{\rm fib},\kappa\varpi_c)$.
\end{corollary}

\begin{proof}
Let $\lambda\in[0,\varpi_{\rm fib})$ satisfy the displayed condition. By
Theorem~\ref{rR6:feshbach}(iii), the definition of $\kappa$ and $\ip g
{A^cg}\ge\varpi_c\lVert g\rVert^2$,
\[
\tilde{\mathcal E}_\lambda(g)\ge\kappa\varpi_c\lVert g\rVert^2-\frac{\lambda
\varpi_{\rm fib}}{\varpi_{\rm fib}-\lambda}\zeta\lVert g\rVert^2\ge\lambda
\lVert g\rVert^2,
\]
and Theorem~\ref{rR6:feshbach}(ii) gives $\varpi'\ge\lambda$. The number
$\lambda=\min(\varpi_{\rm fib}/2,\kappa\varpi_c/(1+2\zeta))$ satisfies the
condition, because $\varpi_{\rm fib}/(\varpi_{\rm fib}-\lambda)\le2$ and
$\lambda(1+2\zeta)\le\kappa\varpi_c$. For the upper bounds, $\varpi'\le
\varpi_{\rm fib}$ by restriction to $\mathcal H_f$, and the function $f=g+
k_0(g)$ has $\ip fAf=\tilde{\mathcal E}_0(g)$ and $\lVert f\rVert^2=\lVert
g\rVert^2+\lVert k_0(g)\rVert^2$. Dropping $\lVert k_0(g)\rVert^2$,
$\tilde{\mathcal E}_0(g)\le\mathrm K\ip g{A^cg}$ and the infimum over $g$
of $\ip g{A^cg}/\lVert g\rVert^2$ is $\varpi_c$. If $\zeta=0$ and $\kappa=
\mathrm K$, the lower bound holds for every $\lambda<\min(\varpi_{\rm fib},
\kappa\varpi_c)$ and the upper bound is $\min(\varpi_{\rm fib},\kappa
\varpi_c)$.
\end{proof}

In the measure-theoretic case $\mathrm K\le\sup_F\mathcal E'_{n'}(F\circ c)/
\mathcal E^c_n(F)$, since $\tilde{\mathcal E}_0(g)\le\ip gAg$, and under
the hypotheses of \cite[Theorem~94(i)]{BBLR26} recalled above this is at
most $(bn/n')^4$ (the assumption $\tilde{\mathcal E}_0(g)=0$ when $\ip g
{A^cg}=0$ then holds); the upper bound of the corollary is then the descent
$\varpi'\le(bn/n')^4\varpi_c$, sharpened by the fluctuation response. The
lower bound is informative as soon as $\kappa>0$ and $\zeta<\infty$; the
three numbers $\varpi_{\rm fib}$, $\kappa$, $\zeta$ are defined for
every measure with $\varpi_{\rm fib}>0$.

\begin{corollary}[a tower]\label{rR6:tower}
Let $\mu_0,\dots,\mu_K$ be the gauge-invariant measures of a tower as in
\cite[Theorem~101]{BBLR26}, with cube sides $n_j$ and $L_j\ge2n_j+4$, let
$\varpi_j$ be the best constant of $\mathrm{BP}_{\mu_j}(n_j)$, and let
$\varpi_{{\rm fib},j}$, $\kappa_j$, $\zeta_j$ be the quantities of the step
$(\mu_j,c_j)$ in the measure-theoretic case, $1\le j\le K$, and assume
$\varpi_{{\rm fib},j}>0$. Put $r_j=\min(1,\kappa_j/(1+2\zeta_j))$ and
$\varphi=\min_j\varpi_{{\rm fib},j}/2$. Then
\[
\varpi_K\ge\min(\varpi_0,\varphi)\prod_{j=1}^Kr_j .
\]
In particular, if $\kappa_j\ge1-\epsilon_j$ and $\zeta_j\le\epsilon_j$ with
$\epsilon_j\in[0,\frac14]$ and $\sum_j\epsilon_j\le E$, then $\varpi_K\ge
e^{-6E}\min(\varpi_0,\varphi)$.
\end{corollary}

\begin{proof}
The coarse constant of the step $(\mu_j,c_j)$ is $\varpi_{j-1}$, so
Corollary~\ref{rR6:recursion} gives $\varpi_j\ge\min(\varphi,r_j
\varpi_{j-1})$, as $r_j\le\kappa_j/(1+2\zeta_j)$. With $m_0=\varpi_0$ and $m_j=\min(\varphi,r_jm_{j-1})$, induction
gives $\varpi_j\ge m_j$, and since $r_j\le1$, $m_K\ge\min(\varpi_0,\varphi)
\prod_jr_j$. Under the last hypothesis $\kappa_j/(1+2\zeta_j)\ge(1-
\epsilon_j)/(1+2\epsilon_j)\ge(1-\epsilon_j)(1-2\epsilon_j)$, and $\log(1-
x)\ge-2x$ for $0\le x\le\frac12$, applied with $x=\epsilon_j$ and $x=2
\epsilon_j$, gives $r_j\ge e^{-6\epsilon_j}$.
\end{proof}

The quantities $\kappa_j$ and $\zeta_j$ enter only through the product of
the $r_j$, so the constant at the finest level is bounded below uniformly
in the number of steps when the defects $1-r_j$ are summable, given the
condition at the coarsest level and a uniform fibre constant. In the other
direction the upper bound of Corollary~\ref{rR6:recursion} gives $\varpi_j
\le\mathrm K_j\varpi_{j-1}$ (under the assumption made for $\mathrm K$),
hence $\varpi_K\le\varpi_0\prod_{j=1}^K\mathrm K_j$; so if, along a
sequence of towers, the upper ratios satisfy $\mathrm K_j\le1$ and $\sum_j(1-\mathrm
K_j)=\infty$, the constants at the finest level tend to zero for bounded
$\varpi_0$. The two products bracket the constant, and they agree, up to the cap by
the fibre constants, when $\zeta_j=0$ and $\kappa_j=\mathrm K_j$. For the
free scalar field below the fluctuation response of the constant linear
function vanishes and its ratio is $1-\delta_j$ with summable $\delta_j$;
this concerns one function, not the infimum $\kappa_j$ or the supremum
$\zeta_j$.

\subsubsection{The constant function and the free field}

For measures on scalar fields with a linear block map the constant linear
function gives an exact recursion with one defect, valid for every
translation-invariant measure. Let $d\ge1$, $b\ge2$, $K\ge1$, $n_0\ge2$,
$L_0\ge n_0+2$, and for $0\le j\le K$ let $L_j=b^jL_0$, $n_j=b^jn_0$ and
$m_j=b^{K-j}$. Let $\nu$ be a probability measure on $\R^{\mathbb T_K}$,
$\mathbb T_K=(\Z/L_K\Z)^d$, invariant under the translations of $\mathbb
T_K$, with $\int\phi_x^2\,d\nu<\infty$ and $\chi=L_K^{-d}\Var_\nu(\sum_x
\phi_x)>0$. The level-$j$ field is $\Phi^{(j)}_y=\sum_{x\in B^{(j)}_y}\phi_x$,
$y\in\mathbb T_j=(\Z/L_j\Z)^d$, with $B^{(j)}_y=m_jy+[0,m_j)^d$, and $\nu_j$
is its law; $\nu_K=\nu$. The cubes of $\mathbb T_j$ have side $n_j$, so that
the cube $V_j=[0,n_j)^d$ of $\mathbb T_j$ is the union of the level-$j$
blocks of the fine cube $V'=[0,n_K)^d$ for every $j$. Let $X=\sum_{x\in V'}
\phi_x=\sum_{y\in V_j}\Phi^{(j)}_y$, $\mathfrak G_j=\sigma(\Phi^{(j)}_y:y
\notin V_j)$ and
\[
\sigma_j^2=\nu\bigl(\Var(X\mid\mathfrak G_j)\bigr),\qquad\delta_j=1-
\frac{\sigma_j^2}{\sigma_{j-1}^2}\quad(1\le j\le K).
\]
Assume $\sigma_K>0$; then, by Proposition~\ref{rR6:zeromode}(ii) below,
every $\sigma_j>0$. For a function $f\in L^2(\nu_j)$ put $\mathcal E^{(j)}
_{n_j}(f)=n_j^{-d}\sum_{y\in\mathbb T_j}\nu_j((f-P^{(j)}_yf)^2)$, with
$P^{(j)}_y$ the conditional expectation given the $\Phi^{(j)}_z$, $z\notin
V_j+y$, the Dirichlet form of~\eqref{eq:ffvarpi} for $\nu_j$. Let $\ell=
\sum_x\phi_x-\nu(\sum_x\phi_x)$, which equals $\sum_y\Phi^{(j)}_y$
minus a constant for every $j$, let $\varpi^{(0)}_j=\mathcal E^{(j)}_{n_j}
(\ell)/\Var(\ell)$ be its Rayleigh quotient, and let $\varpi_j$ be the
best constant of $\mathrm{BP}_{\nu_j}(n_j)$ on $L^2(\nu_j)$, the infimum
of $\mathcal E^{(j)}_{n_j}(f)/\Var(f)$ over nonconstant $f$.
\begin{proposition}[the recursion on the constant function]\label{rR6:zeromode}
In this setting:
\begin{enumerate}[label=\textup{(\roman*)},leftmargin=*]
\item $\varpi_j\le\varpi^{(0)}_j=\sigma_j^2/(n_K^d\chi)$ for $0\le j\le K$;
\item $\sigma_K^2\le\sigma_{K-1}^2\le\dots\le\sigma_0^2$, so $\delta_j\in
[0,1)$, and
\[
\varpi^{(0)}_j=(1-\delta_j)\,\varpi^{(0)}_{j-1},\qquad\prod_{j=1}^K(1-
\delta_j)=\frac{\sigma_K^2}{\sigma_0^2}=\frac{\varpi^{(0)}_K}{\varpi^{(0)}_0};
\]
\item for the step from $\nu_j$ to $\nu_{j-1}$, $1\le j\le K$, in the
measure-theoretic case of Theorem~\textup{\ref{rR6:feshbach}} with
$\mathcal H$ the real functions in $L^2(\nu_j)$ of mean zero, $\Pi_y=
P^{(j)}_y$, $w'=n_j^{-d}$, $\mathfrak B=\sigma(\Phi^{(j-1)})$, $\mathfrak
G_z=\sigma(\Phi^{(j-1)}_{z'}:z'\notin V_{j-1}+z)$ and $w=n_{j-1}^{-d}$
(so that $\ip{F\circ\Phi^{(j-1)}}{A^c\,F\circ\Phi^{(j-1)}}=\mathcal
E^{(j-1)}_{n_{j-1}}(F)$, $\nu_{j-1}$ being the law of the block sums
$\Phi^{(j-1)}$ of $\Phi^{(j)}$), $\tilde{\mathcal E}_0(\ell)\le(1-\delta_j)\,
\mathcal E^{(j-1)}_{n_{j-1}}(\ell)$, so $\kappa_j\le1-\delta_j$ when the
fibre constant of the step is positive.
\end{enumerate}
\end{proposition}

\begin{proof}
(i) $\Var(\ell)=L_K^d\chi$. The translation of $\mathbb T_j$ by $y$ is the
translation of $\mathbb T_K$ by $m_jy$, which preserves $\nu$, permutes the
level-$j$ blocks and fixes $\ell$; it maps the cube $V_j$ to $V_j+y$. Since
$\ell-X$, which is $\sum_{y\notin V_j}\Phi^{(j)}_y$ minus a constant, is
$\mathfrak G_j$-measurable, $\nu(
\Var(\ell\mid\mathfrak G_j))=\sigma_j^2$, and every cube of $\mathbb T_j$
contributes the same. So $\mathcal E^{(j)}_{n_j}(\ell)=n_j^{-d}L_j^d\sigma
_j^2$, and with $L_K=m_jL_j$ and $n_K=m_jn_j$ the quotient is $\sigma_j^2/
(n_K^d\chi)$. It bounds the best constant, $\ell$ being a test function.

(ii) For $z\notin V_{j-1}$ the block $B^{(j-1)}_z$ is disjoint from $V'$,
which is a union of level-$(j-1)$ blocks, and it is a union of level-$j$
blocks $B^{(j)}_y$, all with $y\notin V_j$. So $\Phi^{(j-1)}_z$ is a sum of
generators of $\mathfrak G_j$, and $\mathfrak G_{j-1}\subset\mathfrak G_j$.
Hence $\nu(\Var(X\mid\mathfrak G_j))=\lVert X-\E[X\mid\mathfrak G_j]
\rVert^2\le\lVert X-\E[X\mid\mathfrak G_{j-1}]\rVert^2$. The identity
follows from (i), and the product telescopes.

(iii) $\ell$ is a function of $\Phi^{(j-1)}$ of mean zero, so $\ell\in
\mathcal H_c$ and $\tilde{\mathcal E}_0(\ell)\le\ip\ell{A\ell}=\mathcal
E^{(j)}_{n_j}(\ell)$, and by (i) and (ii) the ratio of $\mathcal E^{(j)}_{n_j}(\ell)$ to
$\mathcal E^{(j-1)}_{n_{j-1}}(\ell)$ is $\varpi^{(0)}_j/\varpi^{(0)}_{j-1}=
1-\delta_j$, as both have the denominator $\Var(\ell)$. As $\ip\ell
{A^c\ell}=\mathcal E^{(j-1)}_{n_{j-1}}(\ell)>0$ by (i), $\kappa_j\le
\tilde{\mathcal E}_0(\ell)/\ip\ell{A^c\ell}$.
\end{proof}

So the Rayleigh quotient of the constant function obeys an exact recursion
whose only defect is the information that the fine exterior of a cube
carries on its mass beyond the block sums outside it. For the free field
this quotient is the exact constant at the finest level, the fluctuation
response of $\ell$ vanishes at every step, and the defects are bounded by
the inverse cube size. Let $\gamma$ be the centred Gaussian measure on
$\R^{\mathbb T_K}$ with precision $-\Delta+\mu^2$, $\mu>0$, as in
Sec.~\ref{ssec:fffree}, and $s=n_K\mu$; then $\chi=1/\mu^2$.

\begin{proposition}[the free field]\label{rR6:gauss}
Let $\nu$ in the setting above be $\gamma$, and put
\[
\Gamma_d(s)=\frac{7d}4\,(4d+s^2)\Bigl(\frac65\Bigr)^{d-1}.
\]
\begin{enumerate}[label=\textup{(\roman*)},leftmargin=*]
\item $\varpi_K=\varpi^{(0)}_K=\varpi^{(d)}_{n_K}(s)$, the exact constant of
Theorem~\textup{\ref{ffs:thm:exact}}, and $\varpi^{(0)}_j=\mu^2\sigma_j^2/
n_K^d$ for every $j$.
\item For every step the fluctuation response of $\ell$ vanishes: $k_0(\ell)
=0$ and $\tilde{\mathcal E}_0(\ell)=\mathcal E^{(j)}_{n_j}(\ell)=(1-\delta_j)
\mathcal E^{(j-1)}_{n_{j-1}}(\ell)$, so $\kappa_j\le1-\delta_j\le\mathrm K_j$.
\item For $0\le j\le K$,
\[
\sigma_j^2-\sigma_K^2\le\frac{7d}{12}\,\frac{n_K^{d+2}}{n_j},\qquad\sigma_K^2
\ge\Bigl(\frac56\Bigr)^{d-1}\frac{n_K^{d+2}}{12d+3s^2};
\]
hence, for $1\le j\le K$,
\[
\delta_j\le\frac{\Gamma_d(s)}{n_{j-1}},\qquad\sum_{j=1}^K\delta_j\le
\frac{b\,\Gamma_d(s)}{(b-1)n_0},
\]
and
\[
\Bigl(1+\frac{\Gamma_d(s)}{n_0}\Bigr)^{-1}\le\prod_{j=1}^K(1-\delta_j)\le1 .
\]
\item $\varpi^{(0)}_0=s^2\sigma_0^2/n_K^{d+2}$, and
\begin{gather*}
\varpi^{(d)}_{n_K}(s)\le\varpi^{(0)}_0\le\Bigl(1+\frac{\Gamma_d(s)}{n_0}
\Bigr)\varpi^{(d)}_{n_K}(s),\\
\varpi^{(0)}_0\le\Bigl(1+\frac{\Gamma_d(s)}{n_0}\Bigr)\frac{s^2(n_K+1)
(n_K+2)}{12n_K^2}.
\end{gather*}
\end{enumerate}
\end{proposition}

\begin{proof}
(i) $\Var(\ell)=\ip{\mathbf 1}{(-\Delta+\mu^2)^{-1}\mathbf 1}=L_K^d/\mu^2$.
Given the field outside $V'$, $X$ has the conditional variance $\sigma_K^2=
\ip{\mathbf 1_{V'}}{(-\Delta_D+\mu^2)^{-1}\mathbf 1_{V'}}$, with $-\Delta_D$
the Dirichlet Laplacian of $V'$ (Lemma~\ref{lem:fflinear}), so
$\varpi^{(0)}_K=\mu^2\sigma_K^2/n_K^d$ is the constant of
Theorem~\ref{ffs:thm:exact}(iii), which applies since $L_K\ge n_K+2$. The
formula for $\varpi^{(0)}_j$ is Proposition~\ref{rR6:zeromode}(i).

(ii) The law $\nu_j$ is Gaussian with a positive definite covariance, the
conditional expectation given a set of coordinates or given the linear
functions $\Phi^{(j-1)}$ maps linear functions to linear functions, and so
$A$ maps the space of linear functions of $\Phi^{(j)}$ into itself. It
commutes with the translations of $\mathbb T_j$, which preserve $\nu_j$ and
permute the cubes, and $\ell$ is translation invariant; so $A\ell$ is a
translation-invariant linear function $\sum_yh_y\Phi^{(j)}_y$, and since the
covariance is nondegenerate, $h$ is constant and $A\ell=c\,\ell$. As $\ell=
\sum_z\Phi^{(j-1)}_z$ is measurable for the block variables, $A_{fc}\ell=
(I-Q)A\ell=0$, so $k_0(\ell)=0$ and $\tilde{\mathcal E}_0(\ell)=\ip\ell
{A\ell}$. The rest is Proposition~\ref{rR6:zeromode}(iii) and the
definition of $\mathrm K_j$ (here $\varpi_{{\rm fib},j}\ge\varpi_j>0$).

(iii) Let $\mathfrak F=\mathfrak G_K$, $j<K$, and $m=m_j\ge2$, $n=n_K$.
Since $\mathfrak G_j\subset\mathfrak F$, $\sigma_j^2-\sigma_K^2=\lVert\E[X
\mid\mathfrak F]-\E[X\mid\mathfrak G_j]\rVert^2\le\lVert\E[X\mid\mathfrak F]
-Y\rVert^2$ for every $\mathfrak G_j$-measurable $Y$. Let $\tau=(-\Delta_D+
\mu^2)^{-1}\mathbf 1_{V'}$. The conditional mean of $\phi_{V'}$ given the
exterior is $(-\Delta_D+\mu^2)^{-1}$ applied to the sum of the exterior
neighbours, so $\E[X\mid\mathfrak F]=\sum_{w}p(w)\phi_w$, the sum over the
sites $w\notin V'$ adjacent to $V'$, with $p(w)=\tau(u_w)$ for the unique
neighbour $u_w\in V'$ of $w$ (unique since $L_K\ge n+2$). If $w=u\pm e_i$,
then $u_i\in\{0,n-1\}$. The function $v(u)=\tau_1(u_i)$, $\tau_1(t)=(t+1)
(n-t)/2$, satisfies $(-\Delta_D+\mu^2)v\ge\mathbf 1_{V'}$, because $\tau_1$
has second difference $-1$ and vanishes at $-1$ and $n$, while the
Dirichlet differences in the other coordinates are $b_k(u)v(u)\ge0$ as in
the proof of Corollary~\ref{ffs:cor:closed}(iv); the Green function of
$-\Delta_D+\mu^2$ is positive (Theorem~\ref{ffs:thm:exact}(ii)), so $\tau
\le v$ and $p(w)\le\tau_1(0)=n/2$.

The sites $w$ lie in the $2dn_j^{d-1}$ face blocks $B^{(j)}_y$ with one
coordinate $y_i\in\{-1,n_j\}$ and the others in $[0,n_j)$; these are
distinct and disjoint from $V'$ because $L_j\ge n_j+2$, and in each the
sites $w$ fill the slice adjacent to $V'$. Put $Y=\sum_wp(w)m^{-d}
\Phi^{(j)}_{y(w)}$, $y(w)$ the block of $w$, which is $\mathfrak
G_j$-measurable. Then $\E[X\mid\mathfrak F]-Y=\ip h\phi$ with $h=\sum_Bh_B$,
$h_B=\sum_{w\in B}p(w)(\delta_w-m^{-d}\mathbf 1_B)$, and $\sum_xh(x)=0$. In
Fourier variables $(-\Delta+\mu^2)^{-1}\le(-\Delta)^{+}$ on functions of sum
zero, and for every flow $\theta$ on the links with net outflow $h$,
$\ip hv=\sum_e\theta_e(v(e^-)-v(e^+))\le\lVert\theta\rVert\lVert\nabla v
\rVert$, so
\[
\Var(\ip h\phi)\le\ip h{(-\Delta)^+h}=\sup_v\frac{\ip hv^2}{\lVert\nabla v
\rVert^2}\le\lVert\theta\rVert^2 .
\]
In a face block write the sites as $(z,t)$, $z\in F=[0,m)^{d-1}$ along the
face and $t\in[0,m)$ the distance from the slice of the sources. The flow
$\theta^1$ that carries $p(z)(1-(t+1)/m)$ from $(z,t)$ to $(z,t+1)$ has net
outflow $\sum_zp(z)(\delta_{(z,0)}-m^{-1}\mathbf 1_{\{z\}\times[0,m)})$ and
energy $\sum_zp(z)^2\sum_{r=1}^{m-1}(r/m)^2\le m^dn^2/12$. The remainder
$\psi(z,t)=(p(z)-\bar p)/m$, $\bar p$ the mean of $p$ over $F$, is constant
in $t$ with mean zero on each slice; in each slice take the flow of least
energy with net outflow $\psi(\cdot,t)$, of energy $\ip{\psi}{(-\Delta^N_F)
^+\psi}\le\lVert\psi(\cdot,t)\rVert^2/\lambda_1$ with $-\Delta^N_F$ the
Neumann Laplacian of $F$ and $\lambda_1=4\sin^2(\pi/(2m))\ge4/m^2$ its
spectral gap. Summed over the $m$ slices this flow $\theta^2$ has energy at
most $m\cdot\frac{m^2}4\cdot m^{-2}\sum_z(p(z)-\bar p)^2\le m^dn^2/16$ (for
$d=1$, $\psi=0$). With $\theta_B=\theta^1+\theta^2$, $\lVert\theta_B
\rVert^2\le2(\frac1{12}+\frac1{16})m^dn^2=\frac7{24}m^dn^2$ (the two flows
use disjoint sets of links, so the factor $2$ could be dropped), and the
flows of distinct blocks use disjoint links. So $\sigma_j^2-\sigma_K^2\le2dn_j^{d-1}
\frac7{24}m^dn^2=\frac{7d}{12}n^{d+2}/n_j$, using $n=mn_j$.

For the lower bound let $\varphi(t)=(t+1)(n-t)$ and $v=\varphi^{\otimes d}$
on $V'$; with $A=\sum_t\varphi(t)=n(n+1)(n+2)/6$, $B=\sum_t\varphi(t)^2=
n(n+1)(n+2)(n^2+2n+2)/30$ and the Dirichlet energy $2A$ of $\varphi$
(Lemma~\ref{lem:ffparabola}), the Cauchy--Schwarz inequality gives
\[
\sigma_K^2\ge\frac{\ip{\mathbf 1}v^2}{\ip v{(-\Delta_D+\mu^2)v}}=\frac{A^{2d}}
{2dAB^{d-1}+\mu^2B^d}\ge\frac{A(A^2/B)^{d-1}}{2d+s^2/2},
\]
since $\mu^2B/A=\mu^2(n^2+2n+2)/5\le s^2/2$ for $n\ge2$. With $A\ge n^3/6$
and $A^2/B=\frac56n(n+1)(n+2)/(n^2+2n+2)\ge\frac56n$ this is the bound. The
two bounds give $(\sigma_{j-1}^2-\sigma_K^2)/\sigma_K^2\le\Gamma_d(s)/
n_{j-1}$, and $\delta_j=(\sigma_{j-1}^2-\sigma_j^2)/\sigma_{j-1}^2\le(\sigma
_{j-1}^2-\sigma_K^2)/\sigma_K^2$. The sum is geometric, and the product is
$\sigma_K^2/\sigma_0^2\ge(1+\Gamma_d(s)/n_0)^{-1}$.

(iv) The formula is (i) with $\mu^2=s^2/n_K^2$; the bounds follow from
$\varpi^{(0)}_0=\varpi^{(0)}_K/\prod_j(1-\delta_j)$ and (iii). Finally
$\sigma_K^2=\ip{\mathbf 1}\tau\le\ip{\mathbf 1}v=n^{d-1}A/2$ by the
comparison above, and $\sigma_0^2\le(1+\Gamma_d(s)/n_0)\sigma_K^2$.
\end{proof}
\begin{remark}[the massless limit and the photon]\label{rR6:massless}
Let $\mu\to0$ at fixed $n_K$, so that $s\to0$. By
Proposition~\ref{rR6:gauss}(iii) the defects stay summable, with
$\Gamma_d(0)=7d^2(6/5)^{d-1}$, and their product stays in
$[(1+\Gamma_d(s)/n_0)^{-1},1]$; on the constant function the steps of the
tower lose a bounded factor uniformly in $K$. The constant nevertheless tends to zero: by (i) and the
two bounds on $\sigma_K^2$ in the proof of (iii) and (iv),
\[
\Bigl(\frac56\Bigr)^{d-1}\frac{s^2}{12d+3s^2}\le\varpi_K=\frac{\mu^2
\sigma_K^2}{n_K^d}\le\frac{s^2(n_K+1)(n_K+2)}{12n_K^2},
\]
so $\varpi_K\asymp s^2$ uniformly in $n_K\ge2$ for $0<s\le1$, in agreement with
$\varpi^{(d)}_\infty(s)=T_ds^2+O(s^4)$ of Theorem~\ref{ffs:thm:limit}
($T_1=1/12$ is the limit of the upper bound for $d=1$). In the recursion
$\varpi_K=\varpi^{(0)}_0\prod_j(1-\delta_j)$ the factor $s^2$ sits in
$\varpi^{(0)}_0=\sigma_0^2/(n_K^d\chi)$: the conditional variance $\sigma_0^2
\le(1+\Gamma_d(s)/n_0)\sigma_K^2$ of the mass of the cube given the block
sums outside it at the coarsest scale stays bounded as $\mu\to0$, while the
susceptibility $\chi=1/\mu^2$ diverges. In this factorisation the massless
failure comes from the coarsest-scale factor, not from an accumulation of
defects. The bound $\varpi\le8\pi^2(n+1)^2/L^2$ of
Theorem~\ref{rD:maxwell} for the photon is of the same order,
$2(n+1)^2\lvert p\rvert^2$ with $\lvert p\rvert=2\pi/L$ the least nonzero
momentum of the torus in place of $\mu$, and its slow mode plays the part
of $\ell$. That mode is not translation invariant, so
Proposition~\ref{rR6:zeromode} does not apply to it as stated, and no
gauge-invariant form of the recursion is claimed here.
\end{remark}

\begin{remark}[numerical checks]\label{rR6:numerics}
The following values were computed in double precision for $b=2$, $L_0=8$,
$n_0=2$ and $d=1$; no statement rests on them. For $K=4$ and $s=1$ the best
constants $\varpi_j$ of $\mathrm{BP}_{\nu_j}(n_j)$, computed as the least
eigenvalue of $n_j^{-1}Q_j^{1/2}S_jQ_j^{1/2}$ (Lemma~\ref{lem:fflinear}),
are $0.08253$, $0.08525$, $0.09196$, $0.10609$, $0.13327$ from the finest
level to the coarsest, and each equals $\varpi^{(0)}_j$ to six digits,
although the precision matrices $Q_j$, $j<K$, have positive off-diagonal
entries; that $\varpi_j=\varpi^{(0)}_j$ for $j<K$ is observed, not proved.
For $K=6$ the defects from the finest step to the coarsest are $0.0088$,
$0.0215$, $0.0442$, $0.0819$, $0.138$, $0.206$ at $s=1$, roughly halving
toward the fine scales as Proposition~\ref{rR6:gauss}(iii) predicts, and
$\prod_j(1-\delta_j)=0.548$, $0.582$, $0.781$ at $s=10^{-4}$, $1$, $4$. The
ratio of $\sigma_j^2-\sigma_K^2$ to the bound of (iii) was at most $0.25$
in all cases checked ($d=1,2$), and the ratio of the lower bound for
$\sigma_K^2$ to $\sigma_K^2$ lay between $0.47$ ($d=2$, $s=3$) and $0.98$
($d=1$, $s=10^{-4}$). On linear functions of the block variables
(the first chaos) the step quantities were $\varpi^{(1)}_{\rm fib}\in
[1.03,1.16]$, $1-\kappa^{(1)}_j=0.052$, $0.116$, $0.204$, $0.292$ for $j=4,
\dots,1$ and $\zeta^{(1)}\le0.003$ for $d=1$, $s=1$, $K=4$; and
$\varpi^{(1)}_{\rm fib}\in[0.87,0.98]$, $\kappa^{(1)}\ge0.66$,
$\zeta^{(1)}\le0.01$ for $d=2$, $s=1$, $K=2$, $L_0=6$ (the measures are
Gaussian, $A$ and $Q$ preserve the first chaos, and the minimiser of a
linear block function is linear). These restrictions are one-sided:
$\varpi^{(1)}_{\rm fib}\ge\varpi_{\rm fib}$, $\kappa^{(1)}\ge\kappa$,
$\zeta^{(1)}\le\zeta$, so they are consistent with the open lemma (L) for
the free field but do not test it.
\end{remark}

\begin{remark}[the convention of a fixed cube side]\label{rR6:convention}
The cube sides $n_j=b^jn_0$ cover a fixed physical cube, as in
\cite[Corollary~102]{BBLR26}, where $n_ka_k=\ell$. If instead the same
side $n$ is used at every level, the cube of side $n$ of $\mathbb T_j$ is
the union of the level-$j$ blocks of a cube of side $N_j=nm_j$ of $\mathbb
T_K$. For the free field Proposition~\ref{rR6:zeromode}(i), with $n$ in
place of $n_j$, bounds the best constant of $\mathrm{BP}_{\nu_j}(n)$ by the
quotient $\mu^2\gamma(\Var(X_{N_j}\mid\mathfrak G_j))/N_j^d$, with
$X_{N_j}$ the mass of the cube of side $N_j$ and $\mathfrak G_j$ the block
sums outside it, and, when $L_j\ge n+2$,
Proposition~\ref{rR6:gauss}(iii), (iv) applied to the tower with levels
$j,\dots,K$ and cube sides $nb^i$ place this quotient between
$\varpi^{(d)}_{N_j}(N_j\mu)$ and $(1+\Gamma_d(N_j\mu)/n)\varpi^{(d)}_{N_j}
(N_j\mu)$. By Theorem~\ref{ffs:thm:limit} this upper bound for the best
constant is close to $T_d(nm_j\mu)^2$ when $nm_j\mu$ is small and $nm_j$
large, so it decreases roughly by the factor $b^{-2}$ at each step toward
the fine scales, the dependence on
$s^2$ of Sec.~\ref{ssec:fffree} read along the tower. The two conventions
are related at the level of $\mathrm{BP}$ constants by the analogue of
Proposition~\ref{prop:bpscales}: part (i) passes from a smaller to a larger
cube with the factor of the fourth power of the ratio of sides (the $d$-th
power for scalar fields), and part (iii) passes back under an interior
condition.
\end{remark}

\subsubsection{The angle between block and fibre functions}

The quantities $\kappa$ and $\zeta$ are controlled by one constant, the
cosine of the angle between $\mathcal H_c$ and $\mathcal H_f$ in the form
$A$, together with the fine form of block functions. In the theory of
multilevel methods this is the constant of the strengthened
Cauchy--Schwarz inequality~\cite{EV91}, which governs two-level
iterations~\cite{XZ02}. Put, with $0/0=0$,
\begin{gather*}
\gamma=\sup_{g\in\mathcal H_c,\,k\in\mathcal H_f}\frac{\lvert\ip g{Ak}
\rvert}{\ip g{Ag}^{1/2}\ip k{Ak}^{1/2}}\in[0,1],\\
\theta=\inf_g\frac{\ip g{Ag}}{\ip g{A^cg}},
\end{gather*}
the infimum over $g\in\mathcal H_c$ with $\ip g{A^cg}>0$. In the
measure-theoretic case $\ip{F\circ c}{A\,F\circ c}=\mathcal E'_{n'}(F\circ
c)$ is the fine Dirichlet form of a function of the block variables, and
$\gamma=0$ exactly when $A$ maps $\mathcal H_c$ into itself.

\begin{proposition}[the defect through the angle]\label{rR6:cbs}
In the setting of Theorem~\textup{\ref{rR6:feshbach}}, for $g\in\mathcal
H_c$,
\begin{gather*}
(1-\gamma^2)\ip g{Ag}\le\tilde{\mathcal E}_0(g)\le\ip g{Ag},\\
\lVert k_0(g)\rVert^2\le\frac{\ip{k_0(g)}{Ak_0(g)}}{\varpi_{\rm fib}}\le
\frac{\gamma^2\ip g{Ag}}{\varpi_{\rm fib}} .
\end{gather*}
Hence $\kappa\ge(1-\gamma^2)\theta$ and $\zeta\le\gamma^2\lVert A_{cc}
\rVert/\varpi_{\rm fib}$.
\end{proposition}

\begin{proof}
Let $a=\ip g{Ag}$ and, for $k\in\mathcal H_f$, $t=\ip k{Ak}$. Then
$\ip{g+k}{A(g+k)}=a+2\ip g{Ak}+t\ge a-2\gamma\sqrt{at}+t\ge(1-\gamma^2)a$,
since $t-2\gamma\sqrt a\sqrt t\ge-\gamma^2a$; the infimum over $k$ gives the
lower bound, and $k=0$ the upper one. The minimiser $k_0=k_0(g)$ satisfies
$\ip k{A(g+k_0)}=0$ for every $k\in\mathcal H_f$; with $k=k_0$, $\ip{k_0}
{Ak_0}=-\ip g{Ak_0}\le\gamma\sqrt a\,\ip{k_0}{Ak_0}^{1/2}$, so $\ip{k_0}{Ak_0}
\le\gamma^2a$, and $\ip{k_0}{Ak_0}\ge\varpi_{\rm fib}\lVert k_0\rVert^2$
because $k_0\in\mathcal H_f$. The consequences follow from the definitions
of $\kappa$, $\theta$, $\zeta$ and $a\le\lVert A_{cc}\rVert\lVert g
\rVert^2$.
\end{proof}

For the free field and $g=\ell$, $A\ell=c\ell\in\mathcal H_c$, so $\ip\ell
{Ak}=0$ for $k\in\mathcal H_f$: the angle vanishes on $\ell$, which is the
mechanism of Proposition~\ref{rR6:gauss}(ii). The fluctuation response also
turns the separate inequalities (A), (B) into a multiplicative step. \cite[Theorem~100(ii)]{BBLR26} gives $1/\varpi'\le C_A+C_B/
\varpi_c$ under (A) and (B); along a tower with $C_{B,j}=1$ this yields only
$1/\varpi_K\le1/\varpi_0+\sum_jC_{A,j}$, which degenerates linearly in the
number of steps, and the passage from (A), (B) to $\mathrm S$ with the
constants $2C_A$, $2C_B$ costs a factor $2$ per step.

\begin{corollary}[separate inequalities with a summable response]
\label{rR6:separate}
In the setting of Corollary~\textup{\ref{rR6:tower}}, suppose that for
$1\le j\le K$ the separate inequalities \textup{(A)} and \textup{(B)} hold
for $(\mu_j,c_j)$ with constants $C_{A,j}\le C_A$ and $C_{B,j}$, that
$\prod_{i=j}^KC_{B,i}\le\Gamma$ for every $j$, with $\Gamma\ge1$, and that
$\sum_j\zeta_j\le Z$. Then
\[
\varpi_K\ge\frac{\min(\varpi_0,1/(2C_A))}{\Gamma e^{2Z}} .
\]
\end{corollary}

\begin{proof}
For $k\in\mathcal H_f$, $\E[k\mid c_j]=0$, so $\lVert k\rVert^2=\mu_j(\Var
(k\mid c_j))\le C_{A,j}\ip kAk$ by (A), and $\varpi_{{\rm fib},j}\ge1/C_A$.
By (B), $\mathcal E^{(j-1)}_{n_{j-1}}(F)\le C_{B,j}\mathcal E^{(j)}_{n_j}(f)$
for every $f$ with $\E[f\mid c_j]=F\circ c_j$, and the infimum over $f$
gives $\kappa_j\ge1/C_{B,j}$. Corollary~\ref{rR6:recursion} gives
$\varpi_j\ge\min(\varphi,\rho_j\varpi_{j-1})$ with $\varphi=1/(2C_A)$ and
$\rho_j=1/(C_{B,j}(1+2\zeta_j))$. With $m_0=\varpi_0$ and $m_j=\min(\varphi,
\rho_jm_{j-1})$, induction gives $\varpi_K\ge m_K$, and
\[
m_K=\min\Bigl(\varpi_0\prod_{l=1}^K\rho_l,\ \min_{1\le i\le K}\varphi
\prod_{l=i+1}^K\rho_l\Bigr).
\]
Every product $\prod_{l=i+1}^K\rho_l$, $0\le i\le K$, is at least
$\Gamma^{-1}e^{-2Z}$, by $1+2\zeta\le e^{2\zeta}$ and $\Gamma\ge1$ for the
empty product.
\end{proof}

So the loss of the separate inequalities is the product of the constants
$C_{B,j}$ and of the factors $1+2\zeta_j$; for the second to be bounded
along a tower the fluctuation responses must be summable, and by
Proposition~\ref{rR6:cbs} it suffices that the squared angles $\gamma_j^2$
are summable, with a bounded fine form of block functions and a fibre
constant bounded below.
\subsubsection{$\SU(2)$ along the tuning: what is controlled and what is
open}

Consider the tower of \cite[Corollary~102]{BBLR26}: $\mu^{(k)}_{k,L}$ is
the Wilson measure at $\beta_k$, $\mu^{(j-1)}_{k,L}=(c_j)_*\mu^{(j)}_{k,L}$,
$n_j=2^jn_0$, and the quantities $\varpi_{{\rm fib},j}$, $\kappa_j$, $\mathrm
K_j$, $\zeta_j$, $\gamma_j$ are those of the step $(\mu^{(j)}_{k,L},c_j)$ in
the measure-theoretic case of Theorem~\ref{rR6:feshbach}. The lower
bound of Corollary~\ref{rR6:tower} is uniform in $k$ and $L$ when, at every
step, the fibre constant is bounded below, the defect of the
renormalisation ratio $\kappa_j$ is summable and the fluctuation response
is summable; these are sufficient, not shown to be necessary. In the other
direction, by the upper bound of Corollary~\ref{rR6:recursion}, if
$\mathrm K_j\le1$ for all $j$ then the defects $1-\mathrm K_j$ must be
summable for a lower bound uniform in $k$ with $\varpi_0$ bounded. Ba{\l}aban's renormalisation group bears on
these quantities as follows.

\emph{The propagators.} In~\cite{Bal84a} the renormalisation transformations and their
propagators are studied for the vector field with Gaussian action, in~\cite{Bal85a} the averaging operations and their
regularity classes, in~\cite{Bal85c} the existence and uniqueness up to gauge of the minimiser
of the action at given averages (the background field), and
in~\cite{Bal85b} regularity and decay of the propagators that depend on
such a background field. The small-field fluctuation integrals
of~\cite{Bal87,Bal88} are expansions around the background field built on
these propagators. The propagator bounds would be the natural starting
point of a proof of a lower bound for $\varpi_{{\rm fib},j}$ on small
fields, and their decay is the kind of input that would localise the
fluctuation response $k_0$ near the boundary of the cube, as the flows of
the proof of Proposition~\ref{rR6:gauss}(iii) localise the free-field
defect in a layer of blocks. They are statements about covariances and
Green functions, not about the cube-chain form $A$, so they bound neither
$\gamma_j$ nor $\zeta_j$ as they stand.

\emph{The effective actions.} In~\cite{Bal87} a sequence of localised
effective actions is constructed in the small-field approximation by
cluster expansions in one-step transformations, with $\beta$-functions and
a coupling constant renormalisation; in~\cite{Bal88} the fluctuation
integral is represented by an exponentiated cluster expansion, which
completes the construction of the small-field effective actions. These describe, on small fields, the
marginal law $\mu^{(j-1)}_{k,L}$ of the block variables, and with it the
coarse form $\mathcal E^{(j-1)}_{n_{j-1}}$, together with the flow of the
coupling that defines the tuning. The ratio $\kappa_j$ is not a property of
the marginal law: $\tilde{\mathcal E}_0(F\circ c_j)$ is the least fine
Dirichlet form of a function with conditional mean $F\circ c_j$, and for
the analogue of $\ell$ it is, by Proposition~\ref{rR6:zeromode}, the
comparison of a conditional variance given the fine exterior of a cube with
the conditional variance given the block variables outside it. That
comparison concerns the joint law of fine variables near the cube and block
variables away from it, and we are not aware of an estimate
in~\cite{Bal87,Bal88} stated for it.

\emph{The large fields.} In~\cite{Bal89a,Bal89b} the renormalisation
operation $\mathbf R$ of the large-field regions is constructed, localised,
exponentiated and bounded, which completes Ba{\l}aban's proof of
ultraviolet stability of the four-dimensional pure gauge theory: bounds on
the partition function in a box of fixed physical size, uniform in the
spacing. These are upper bounds on integrals over the large-field
configurations. A lower bound for $\varpi_{{\rm fib},j}$ is a Poincar\'e
inequality for the conditional law given the block variables at every
block configuration of positive probability, including the large ones, and
such an inequality is not a consequence of bounds on partition functions.

\emph{The volume.} The ultraviolet-stability bounds of~\cite{Bal89a,Bal89b}
concern a box
of fixed physical size, while the quantities above are needed uniformly in
$L\in\mathcal L_k$, an infinite family of tori of growing physical size.

So on the set of small fields the propagator bounds of~\cite{Bal84a,Bal85b}
are candidates for controlling the Gaussian part of the defect,
through an argument of the kind of Proposition~\ref{rR6:gauss}(iii) with the
background propagator in place of $(-\Delta+\mu^2)^{-1}$; that argument is
not carried out here. We know of no proved estimate that controls the non-Gaussian part of
$\kappa_j$, the fluctuation response on large fields, or the fibre constant
on large fields.

\begin{remark}[why the defect should be suppressed by the cube size]
\label{rR6:af}
Along the asymptotic-freedom tuning the effective coupling of the measure
at level $j$ is, by perturbative running, of order $g_j^2\asymp1/j$ for $j$ large, so $\sum_{j\le k}g_j^2$
grows like $\log k$. If the upper ratios satisfied $\mathrm K_j\le1-cg_j^2$
at every step, the bound $\varpi_k\le\varpi_0\prod_j\mathrm K_j$ that
follows from Corollary~\ref{rR6:recursion} would decay like a negative power of
$k$, that is of $\log(1/a_k)$, and the constant would tend to zero for
bounded $\varpi_0$; a defect $1-\kappa_j\asymp g_j^2$ would make the lower
bound of Corollary~\ref{rR6:tower} degenerate in the same way. In the
convention of a fixed physical cube the weakly coupled steps $j\approx k$ have the largest cubes
in lattice units, $n_{j-1}=2^{j-1}n_0$, and for the free field the defect is
at most $\Gamma_d(s)/n_{j-1}$ (Proposition~\ref{rR6:gauss}(iii)). A uniform
bound therefore needs the per-step defect of the interacting measure to be
suppressed by the cube size, as a boundary-layer effect, and not merely by
the coupling. This remark is heuristic and is used nowhere.
\end{remark}

\begin{definition}[the open lemma (L): summable renormalisation defect]
\label{rR6:lemmaL}
In the setting of \cite[Corollary~102]{BBLR26}, the \emph{open lemma
\textup{(L)}} is the statement that there are $\varphi>0$ and $E<\infty$
such that for every $k$, every $L\in\mathcal L_k$ and every $1\le j\le k$
there is $\epsilon_j=\epsilon_j(k,L)\in[0,\frac14]$ with $\sum_{j=1}^k
\epsilon_j\le E$ and the following hold for $\mu'=\mu^{(j)}_{k,L}$,
$\mu^c=\mu^{(j-1)}_{k,L}$ and $c=c_j$, with $\mathcal E^{(i)}_{n_i}$ the
Dirichlet form of $\mu^{(i)}_{k,L}$ for the cubes of side $n_i$:
\begin{enumerate}[label=\textup{(L\arabic*)},leftmargin=*]
\item (fibre) $2\varphi\,\mu'(\Var(f\mid c))\le\mathcal E^{(j)}_{n_j}(f)$
for every real gauge-invariant $f\in L^2(\mu')$ with $\E[f\mid c]=0$;
\item (renormalisation) $(1-\epsilon_j)\,\mathcal E^{(j-1)}_{n_{j-1}}(F)
\le\mathcal E^{(j)}_{n_j}(f)$ for every real gauge-invariant $F\in L^2
(\mu^c)$ of mean zero and every real gauge-invariant $f\in L^2(\mu')$ with
$\E[f\mid c]=F\circ c$;
\item (response) for every such $F$, the minimiser $f_F$ of $\mathcal
E^{(j)}_{n_j}$, which exists by (L1) and Theorem~\ref{rR6:feshbach}(i),
under the constraint $\E[f\mid c]=F\circ c$ satisfies
$\mu'(\Var(f_F\mid c))\le\epsilon_j\Var_{\mu^c}(F)$.
\end{enumerate}
In the notation of Corollary~\ref{rR6:recursion}: $\varpi_{{\rm fib},j}\ge
2\varphi$, $\kappa_j\ge1-\epsilon_j$ and $\zeta_j\le\epsilon_j$.
\end{definition}

The open lemma (L) is not proved, for $\SU(2)$ or for any interacting
measure; for the free field, restricted to the constant function, (L2) holds with
$\epsilon_j=\delta_j\le\Gamma_d(s)/n_{j-1}$
(Proposition~\ref{rR6:gauss}(iii)), summable, and within the cap
$\epsilon_j\le\frac14$ for the steps with $n_{j-1}\ge4\Gamma_d(s)$, and
(L3) holds since the response vanishes; (L) is not proved on all
functions. Part
(L1) follows from (H2) through (A), with $2\varphi=1/C_A$, and (L2) is (B) with $C_{B,j}=1/(1-
\epsilon_j)$; (L) asks in addition that these constants have summable
defects and that the responses are summable, and (H2) does not ask for
the responses. Neither (L) nor (H2) is shown here to imply the other. The
inequality $\inf\{\mathcal E^{(j)}_{n_j}(f)-\lambda\mu'(\Var(f\mid c)):\E[f
\mid c]=F\circ c\}\ge\rho\,\mathcal E^{(j-1)}_{n_{j-1}}(F)$ for all $F$, the
form of (L2) at a spectral parameter $\lambda>0$, is exactly $\mathrm S(n_j,
n_{j-1};1/\lambda,1/\rho)$, since $\mathcal E^{(j)}_{n_j}(f)-\lambda\mu'
(\Var(f\mid c))\ge\rho\,\mathcal E^{(j-1)}_{n_{j-1}}(F_f)$ for every $f$ is
the split inequality; (L2) is its form at $\lambda=0$, whose analogue for
scalar fields has on the constant function the defect $\delta_j$ of
Proposition~\ref{rR6:zeromode}.

\begin{proposition}[the conditional consequence]\label{rR6:conditional}
In the setting of \cite[Corollary~102]{BBLR26}, suppose that
\textup{(H1)} and the open lemma \textup{(L)} hold. Then $\mathrm{BP}(n_k,
\varpi)$ holds for $\mu_{k,L}$ for all $k$ and $L\in\mathcal L_k$, with
\[
\varpi=e^{-6E}\min(\varpi_0,\varphi),
\]
and every continuum limit as in Theorem~\textup{\ref{thm:avgtransfer}}
satisfies $\sigma(H)\subset\{0\}\cup[16\varpi/(625e^2\ell),\infty)$.
\end{proposition}

\begin{proof}
Fix $k$ and $L$. By (H1) the best constant of $\mathrm{BP}$ for
$\mu^{(0)}_{k,L}$ at the side $n_0$ is at least $\varpi_0$. By (L1) the
number $\min_j\varpi_{{\rm fib},j}/2$ of Corollary~\ref{rR6:tower} is at
least $\varphi$, and by (L2), (L3) $\kappa_j\ge1-\epsilon_j$ and $\zeta_j\le
\epsilon_j$ with $\epsilon_j\in[0,\frac14]$ and $\sum_j\epsilon_j\le E$.
Corollary~\ref{rR6:tower}, whose tower hypotheses hold since $L_j\ge2n_j+4$
as in the proof of \cite[Corollary~102]{BBLR26}, gives that the best
constant of $\mathrm{BP}(n_k)$ for $\mu_{k,L}$ is at least $e^{-6E}\min(
\varpi_0,\varphi)$. Theorem~\ref{thm:avgtransfer} applies with $n_ka_k=
\ell$.
\end{proof}

The proposition is conditional. By Theorem~\ref{thm:avgtransfer} its
conclusion gives a mass gap only for a nontrivial continuum limit, and
neither (H1), nor (L), nor the nontriviality of the limit is proved;
(D) and (E) are not addressed.

\begin{remark}[what is proved and what is open]\label{rR6:status}
Proved, for every probability measure and every sub-$\sigma$-algebra of
block variables with a positive fibre constant: the exact characterisation of the fine constant by the
effective forms (Theorem~\ref{rR6:feshbach}); the two-sided recursion in
terms of $\varpi_{\rm fib}$, $\kappa$, $\zeta$ and $\mathrm K$
(Corollary~\ref{rR6:recursion}) and its tower form
(Corollary~\ref{rR6:tower}); the control of $\kappa$ and $\zeta$ by the
angle $\gamma$ (Proposition~\ref{rR6:cbs}); the multiplicative step from
the separate inequalities (Corollary~\ref{rR6:separate}); and, for
translation-invariant measures with block sums, the exact recursion on the
constant function (Proposition~\ref{rR6:zeromode}). For the massive free
field: the recursion reproduces the exact constant of
Theorem~\ref{ffs:thm:exact}, the response of $\ell$ vanishes, and the
defects are at most $\Gamma_d(s)/n_{j-1}$, so their product lies in
$[(1+\Gamma_d(s)/n_0)^{-1},1]$ (Proposition~\ref{rR6:gauss}); as $\mu\to0$
the constant is of order $s^2$, and in the factorisation the factor $s^2$
sits in the coarsest-scale quotient (Remark~\ref{rR6:massless}). Open: whether $\varpi_j=\varpi^{(0)}_j$ for
$j<K$ in the free field (observed numerically); the open lemma (L) on all
functions, already for the free field; every quantity of the step for
$\SU(2)$, in particular (L1) on large fields, the non-Gaussian part of
(L2) and (L3); a gauge-invariant form of the recursion for the photon; and
(H1).
\end{remark}
